\documentclass[11pt,a4paper]{article}
\usepackage[margin=27mm]{geometry}
\usepackage{amsmath,amssymb,fontspec}
\setmainfont{STIXGeneral}
\usepackage{hyperref}
\title{Independent review of the Topic B proof dependencies}
\author{Independent review agent}
\date{11 October 2026}
\begin{document}
\maketitle

\section{What was reviewed}
The review concerns the mathematical meaning of the paired dependency graph,
not a new proof of the historical classification inputs and not a measurement
of compression. I read the complete current common paper and its four
applications, and compared the historical routes with the relevant complete
proof bodies of N001, N003, N008 and the fixed revised edition of L030.
For N001 the available document is the previously bound authored browser
rendering. Its original PDF remains unavailable. The graph records thirty
fixed clauses: four, eight, eleven and seven respectively.

The principal organizing proofs in the common paper are correctly separated.
The forest intersection argument determines the symmetric subgroups; the
laminar partition argument supplies a common representative for a flag.
Neither conclusion is substituted for the other. Relative generation uses
the successive maximal intersections and the panel choices. The first
application retains the earlier classification and uniqueness of inductively
minimal pairs. The second retains the stated Whiston inputs and their
classification dependence, then supplies the tree and parity calculations.
The third proves the necessary finite relative-generation step before using
the Frattini ladder, including for an infinite ambient group. The fourth
proves its incidence, distance and automorphism assertions, with a separate
spectral argument for automorphisms of the uncolored graph.

All arrows were assessed as prerequisites pointing to uses. A displayed
route need not be a minimal proof. The parent and branch locators are
semantic references, not disjoint quantities that can be added as costs.
Attribution, formal verification and diagnostic evidence remain separate
from the mathematical arrows. No required enumeration program is hidden
in this topic. Retained external costs and unavailable image costs remain
unknown, not zero.

\section{Five corrections to the graph}
The initial frozen graph had no newly discovered manuscript proof defect,
but its finer routes needed the following repairs. The revised graph and
its exact hashes are bound by the accompanying review record.
\begin{enumerate}
\item The propagation branch in N008, Theorem 3.1, uses the maximal chain,
the first failed index and earlier containment hypotheses established in
the preceding setup. An arrow from that setup was missing. The final
contradiction belongs to the assembled theorem; the setup locator is not a
claim that the propagation proof depends on its own conclusion.
\item The replacement route for N003, Corollary 2.3, needs the classification
of an arbitrary extremal independent family before applying the constructive
forest and parity intersection formulas. The full application already
contained this classification; the individual output route omitted it.
\item Clause C11 includes containment of the Frattini subgroup in the action
kernel for a firm residually weakly primitive geometry. The general identity
for the core alone does not establish that containment. The corrected route
also includes the Frattini containment theorem. Normality then completes
the implication.
\item The retained classification dependence applies to the cited Whiston
bound as well as to the omission-subgroup trichotomy. The graph must not
make the former appear to use a separately established classification-free
proof. This records the retained proof, not logical indispensability of the
classification theorem.
\item L030's one-line proof of connectedness of the bipartite double mentions
odd cycles but also needs connectedness of the original Kneser graph. The
even and odd paths in Lemma 2.2 supply that premise, including the parameter
boundary not covered by the strict inequality in the printed Lemma 2.1.
The added edge is explicitly an implicit premise, not a fabricated citation.
\end{enumerate}

\section{Scope and limitations}
The graph preserves the already accepted qualifications. The false printed
converse in N003's recall of Whiston is not retained. Its final ``immediate''
corollary is distinguished from the explicit proof for arbitrary parabolic
subgroups in the application. N008's unrestricted infinite generating-family
converse is corrected, while the infinite geometric containment theorem is
preserved. Its dimension exercise is distinguished from the new written
calculation. The L030 corrections concern the exact revised edition and the
displayed 63-vertex graph; they do not refute the existence of a different
locally Desargues graph.

This is a mathematical and documentary graph review. It does not certify
every sentence of the historical papers, reread the original proofs of all
retained classifications, or repeat the previously reviewed Lean and finite
diagnostics. The structural replay checks bindings, spans, acyclicity and
closures. A separate small independent checker protects the five repaired
interfaces and recomputes the thirty paired output closures and the four
complete application closures. Neither checker replaces the written reading.
The acceptance record assigns no text weights, factor, historical priority,
new source coverage or new conjecture credit.
\end{document}
